\documentclass[11pt]{article}
\usepackage[T1]{fontenc}
\usepackage{lmodern}
\input{glyphtounicode}
\pdfgentounicode=1
\usepackage[utf8]{inputenc}
\usepackage[margin=1in]{geometry}
\usepackage{amsmath,amssymb,amsthm,booktabs,hyperref}
\hypersetup{colorlinks=true,urlcolor=blue,pdftitle={Equal-length Tamari orbits with unequal down-degree totals},pdfauthor={Wentao Li},pdfsubject={Mathematical research draft prepared with assistance from OpenAI Codex}}
\newtheorem{theorem}{Theorem}
\newtheorem{lemma}[theorem]{Lemma}
\title{Equal-length Tamari orbits with unequal down-degree totals}
\author{Wentao Li}
\date{November 20, 2026}
\begin{document}\maketitle
For a finite lattice the down-degree of an element is the number of its lower covers. Defant and Lin conjectured that on each $\nu$-Tamari lattice any two rowmotion orbits of equal length have equal total down-degree (homometry). This is weaker than requiring a common average for all orbits.
\begin{theorem}On the ordinary $\nu$-Tamari lattice for $\nu=EN^3E^3NEN$, two orbits have length 22 and down-degree sums 46 and 45. The same failure holds for every alt increment vector for this path.\end{theorem}
The border path $\nu$ bounds the partition $\lambda=(5,4,1,1,1)$. First take the distributive lattice of subdiagrams of $\lambda$, recording an ideal by row lengths $u=(u_1,\ldots,u_5)$ from the top. Its down-degree $d(u)$ is the number of distinct positive row lengths, because one last cell can be removed at each such length. The minimal missing cells are $(i,u_i+1)$ for $u_i<\lambda_i$ and $i=1$ or $u_{i-1}>u_i$ (with the inequality $u_i<\lambda_i$ required in all cases). Let $A(u)$ denote that set. Rowmotion generates the ideal below $A(u)$, so
\[
 (Ru)_i=\max\big(\{j:(h,j)\in A(u),\ h\ge i\}\cup\{0\}\big).
\]
The following are cyclic lists. Each string represents the five row lengths.
\begin{verbatim}
11000 21100 32110 43111 54000 11100 21110 32111 43000 54100 11110
21111 32000 43100 54110 11111 20000 31000 42100 53110 44111 50000

22100 31110 42111 53000 44100 51110 22111 30000 41000 52100 33110
41111 52000 33100 41110 52111 33000 41100 52110 33111 40000 51000
\end{verbatim}
Applying the displayed rule gives the next entry, including wrap-around. Each list has 22 distinct entries and they are disjoint. Summing $d(u)$ gives respectively 46 and 45. In the path convention, nesting can reverse the subdiagram order. Duality preserves these totals: it reverses rowmotion up to the natural identification and interchanges up- and down-degree, which have equal orbit sums because the minimal missing cells become the maximal cells of the next ideal.

Theorem 1.1 of \href{https://arxiv.org/abs/2609.12983}{Adenbaum, Barnard, Ceballos, Chenevière, Defant, Hopkins, Müller, Rubey and Striker} gives an orbit-length and down-degree-sum preserving bijection between the distributive nesting lattice and each alt $\nu$-Tamari lattice, including ordinary $\mathrm{Tam}(\nu)$. Thus the witness disproves Defant--Lin's unrestricted conjecture and the homometry assertion of \href{https://arxiv.org/abs/2605.29431}{Eu, Hioe and Lee, Conjecture 5.1}. Their conjecture also includes independence of the increment vector, which the transfer theorem establishes.

\paragraph{Direct check in the ordinary lattice.} The included program independently builds all 62 paths above $\nu$, computes the ordinary valley rotations and labels each cover $x\lessdot y$ by the minimum of $\{z:z\le y,\ z\not\le x\}$. For a cover this equals the minimum of $\{z:z\vee x=y\}$. Semidistributive rowmotion sends the up-label set at an element to the equal down-label set at its image. The full list of (length, down-degree sum) is $(6,12),(12,25),(22,45),(22,46)$, each once.

\paragraph{Original question.} Colin Defant and James Lin (2022), Conjecture 10.2. \url{https://arxiv.org/abs/2208.10464}.
\medskip\noindent\textit{Research draft prepared with assistance from OpenAI Codex. Independent review pending.}\end{document}
